% =====================================================================
\begin{thebibliography}{99}
\addcontentsline{toc}{chapter}{References}
\small

\bibitem{tan2023} H.~Tan, W.~Zheng and P.~Vijayakumar, ``Secure and efficient authenticated key management scheme for UAV-assisted infrastructure-less
IoVs,'' \emph{IEEE Transactions on Intelligent Transportation Systems}, vol.~24, no.~6, pp.~6389--6400, June 2023.

\bibitem{alriyami} S.~S. Al-Riyami and K.~G. Paterson, ``Certificateless public key cryptography,'' in \emph{ASIACRYPT 2003}, LNCS 2894, Springer,
pp.~452--473, 2003.

\bibitem{gg09} D.~Galindo and F.~D. Garcia, ``A Schnorr-like lightweight identity-based signature scheme,'' in \emph{AFRICACRYPT 2009}, LNCS 5580,
Springer, pp.~135--148, 2009.

\bibitem{ckk12} S.~Chatterjee, C.~Kamath and V.~Kumar, ``Galindo--Garcia identity-based signature revisited,'' in \emph{ICISC 2012}, LNCS 7839,
Springer, pp.~456--471, 2013.

\bibitem{schnorr} C.~P. Schnorr, ``Efficient signature generation by smart cards,'' \emph{Journal of Cryptology}, vol.~4, no.~3, pp.~161--174, 1991.

\bibitem{forking} D.~Pointcheval and J.~Stern, ``Security arguments for digital signatures and blind signatures,'' \emph{Journal of Cryptology},
vol.~13, no.~3, pp.~361--396, 2000.

\bibitem{bnforking} M.~Bellare and G.~Neven, ``Multi-signatures in the plain public-key model and a general forking lemma,'' in \emph{ACM CCS 2006},
pp.~390--399, 2006.

\bibitem{bgr} M.~Bellare, J.~A. Garay and T.~Rabin, ``Fast batch verification for modular exponentiation and digital signatures,'' in
\emph{EUROCRYPT 1998}, LNCS 1403, Springer, pp.~236--250, 1998.

\bibitem{girault} M.~Girault, ``Self-certified public keys,'' in \emph{EUROCRYPT 1991}, LNCS 547, Springer, pp.~490--497, 1991.

\bibitem{chiou89} G.-H. Chiou and W.-T. Chen, ``Secure broadcasting using the secure lock,'' \emph{IEEE Transactions on Software Engineering},
vol.~15, no.~8, pp.~929--934, 1989.

\bibitem{sp24} P.~Sharma and B.~R. Purushothama, ``Proving the (in)security of CRT based key management schemes under SAOA model,''
\emph{Wireless Personal Communications}, vol.~134, no.~3, pp.~1299--1321, 2024.

\bibitem{agh09} M.~Albrecht, C.~Gentry, S.~Halevi and J.~Katz, ``Attacking cryptographic schemes based on `perturbation polynomials','' in
\emph{ACM CCS 2009}, 2009.

\bibitem{wgl98} C.~K. Wong, M.~Gouda and S.~S. Lam, ``Secure group communications using key graphs,'' in \emph{ACM SIGCOMM 1998}, pp.~68--79, 1998.

\bibitem{lll} A.~K. Lenstra, H.~W. Lenstra Jr. and L.~Lov\'asz, ``Factoring polynomials with rational coefficients,'' \emph{Mathematische Annalen},
vol.~261, pp.~515--534, 1982.

\bibitem{cz} D.~G. Cantor and H.~Zassenhaus, ``A new algorithm for factoring polynomials over finite fields,'' \emph{Mathematics of Computation},
vol.~36, no.~154, pp.~587--592, 1981.

\bibitem{hkdf} H.~Krawczyk, ``Cryptographic extraction and key derivation: the HKDF scheme,'' in \emph{CRYPTO 2010}, LNCS 6223, Springer,
pp.~631--648, 2010.

\bibitem{aead} P.~Rogaway, ``Authenticated-encryption with associated-data,'' in \emph{ACM CCS 2002}, pp.~98--107, 2002.

\bibitem{kerberos94} B.~C. Neuman and T.~Ts'o, ``Kerberos: an authentication service for computer networks,'' \emph{IEEE Communications Magazine},
vol.~32, no.~9, pp.~33--38, 1994.

\bibitem{castro99} M.~Castro and B.~Liskov, ``Practical Byzantine fault tolerance,'' in \emph{OSDI 1999}, pp.~173--186, 1999.

\bibitem{bdtw01} D.~Boneh, X.~Ding, G.~Tsudik and C.~M. Wong, ``A method for fast revocation of public key certificates and security capabilities,''
in \emph{10th USENIX Security Symposium}, 2001.

\bibitem{cbn06} S.~S.~M. Chow, C.~Boyd and J.~M. Gonz\'alez Nieto, ``Security-mediated certificateless cryptography,'' in \emph{PKC 2006}, LNCS 3958,
Springer, pp.~508--524, 2006.

\bibitem{yap07} W.-S. Yap, S.~S.~M. Chow, S.-H. Heng and B.-M. Goi, ``Security mediated certificateless signatures,'' in \emph{ACNS 2007}, LNCS 4521,
Springer, 2007.

\bibitem{frost20} C.~Komlo and I.~Goldberg, ``FROST: Flexible round-optimized Schnorr threshold signatures,'' in \emph{Selected Areas in Cryptography
(SAC 2020)}, LNCS 12804, Springer, 2021.

\bibitem{rfc9591} D.~Connolly, C.~Komlo, I.~Goldberg and C.~A. Wood, ``The Flexible Round-Optimized Schnorr Threshold (FROST) protocol for two-round
Schnorr signatures,'' RFC 9591, IRTF, June 2024.

\bibitem{ros21} F.~Benhamouda, T.~Lepoint, J.~Loss, M.~Orr\`u and M.~Raykova, ``On the (in)security of ROS,'' in \emph{EUROCRYPT 2021}; also
\emph{Journal of Cryptology}, vol.~35, article 25, 2022.

\bibitem{dolevyao} D.~Dolev and A.~C. Yao, ``On the security of public key protocols,'' \emph{IEEE Transactions on Information Theory}, vol.~29,
no.~2, pp.~198--208, 1983.

\bibitem{blanchet16} B.~Blanchet, ``Modeling and verifying security protocols with the applied pi calculus and ProVerif,'' \emph{Foundations and
Trends in Privacy and Security}, vol.~1, no.~1--2, pp.~1--135, 2016.

\bibitem{paar} C.~Paar and J.~Pelzl, \emph{Understanding Cryptography}, Springer, 2010 (the curve $y^2=x^3+2x+2$ over $\mathbb F_{17}$ is its running
example in Chapter~9).

\end{thebibliography}
